\section{推理服务：让模型又快又省地回答请求}
\label{sec:infer}

\begin{guide}
\begin{minipage}[c]{0.66\linewidth}
\noindent 模型训练好之后，要面对的是真实用户：首字要出得快，输出要流畅，而显卡数量有限，还要尽量服务更多的人。从哪里入手？本节先把推理拆成预填充与解码两个阶段分别分析，再看 KV 缓存如何避免重复计算，以及 FlashAttention、PagedAttention、连续批处理和量化各自缓解哪一个瓶颈。效果用 TTFT、TPOT、吞吐和尾延迟等指标衡量；各项优化作用于不同环节，可以叠加使用，但彼此不能替代。
\end{minipage}\hfill
\begin{minipage}[c]{0.3\linewidth}\centering\locator{2}\end{minipage}
\end{guide}

\subsection{推理的两个阶段：预填充与解码}
\label{subsec:prefill-decode}

大模型以\emph{自回归}方式生成文本：每次只预测下一个 token，并把它追加到输入末尾，再预测下一个。据此，一次请求的处理可以分为两个阶段（图\,\ref{fig:prefill}）。

\paragraph{预填充（Prefill）} 模型一次性并行处理用户输入的全部 token，计算每一层的注意力 Key 与 Value 并保存下来，同时得到第一个输出 token。由于大量 token 共享同一次权重读取，这一阶段算术强度高，通常\emph{计算受限}。衡量它的指标是\term{首 token 延迟}（\en{Time To First Token}，TTFT），即从发出请求到看到第一个输出 token 的时间，其中还包括请求排队的时间。

\paragraph{解码（Decode）} 模型逐个生成后续 token。每一步只处理一个新 token，却要读取全部权重和此前积累的 Key、Value，因此算术强度很低，通常\emph{访存受限}（见第\,\ref{subsec:roofline}\,小节）。衡量它的指标是\term{每 token 输出时间}（\en{Time Per Output Token}，TPOT），或其倒数——输出速度（token/s）。

\begin{figure}[htbp]
\centering
\begin{tikzpicture}[fig, x=1mm, y=1mm]
  \tikzset{tok/.style={draw=figline, fill=white, minimum size=4.2mm, inner sep=0pt, font=\figfont\footnotesize},
           newtok/.style={draw=accent, fill=accentpale, minimum size=4.2mm, inner sep=0pt, font=\figfont\footnotesize}}
  % ---------- 计算行：整体以 x=0 为中心，范围约 -55 ~ 55 ----------
  \node[blk, minimum width=38mm, minimum height=9mm] (pf) at (-36,0) {预填充\\[-1pt]\footnotesize 并行处理全部输入};
  \foreach \i in {1,...,5} {
    \node[tok] (x\i) at (-57+\i*5,14) {x\textsubscript{\i}};
    \draw[->, line width=0.4pt] (x\i.south) -- (x\i.south |- pf.north);
  }
  \node[newtok] (y1) at (-21,14) {y\textsubscript{1}};
  \draw[->, line width=0.4pt] (y1.south |- pf.north) -- (y1.south);
  \foreach \k [evaluate=\k as \n using int(\k+1)] in {1,...,5} {
    \pgfmathsetmacro{\xx}{-6+(\k-1)*12}
    \node[blk, minimum width=9mm, minimum height=9mm, inner sep=1pt, font=\figfont\scriptsize] (dc\k) at (\xx,0) {解码};
    \node[newtok] (y\n) at (\xx,14) {y\textsubscript{\n}};
    \draw[->, line width=0.4pt] (dc\k.north) -- (y\n.south);
  }
  \draw[->, line width=0.4pt] (pf.east) -- (dc1.west);
  \foreach \k [evaluate=\k as \n using int(\k+1)] in {1,...,4} \draw[->, line width=0.4pt] (dc\k.east) -- (dc\n.west);
  \node[font=\figfont\small] at (51.5,0) {…};
  \node[note, anchor=south] at (-36,17.5) {输入提示（T 个 token）};
  \node[note, anchor=south] at (18,17.5) {逐个生成输出 token};
  % ---------- KV 缓存：居中标注，分段表示逐步增长 ----------
  \node[blkw, minimum width=110mm, minimum height=7mm] (kv) at (0,-16) {KV 缓存（提示部分一次写入，生成时每步追加一项）};
  \draw[->, accent, line width=0.7pt] (pf.south) -- (pf.south |- kv.north);
  \foreach \k in {1,...,5} \draw[<->, warn, line width=0.7pt] (dc\k.south) -- (dc\k.south |- kv.north);
  % ---------- 时间轴与指标 ----------
  \draw[->, line width=0.5pt, rulegray] (-55,-25) -- (57,-25) node[right, note, text=ink] {时间};
  \draw[decorate, decoration={brace, mirror, amplitude=3pt}, line width=0.6pt] (-55,-27) -- (-17,-27)
    node[midway, below=3pt, note] {TTFT：首 token 延迟\\[-1pt]（计算受限）};
  \draw[decorate, decoration={brace, mirror, amplitude=3pt}, line width=0.6pt] (-10.5,-27) -- (-1.5,-27)
    node[midway, below=3pt, note] {TPOT\\[-1pt]（访存受限）};
  % ---------- 居中图例 ----------
  \begin{scope}[shift={(0,-44)}]
    \draw[->, accent, line width=0.7pt] (-44,0) -- (-38,0); \node[note, anchor=west] at (-37,0) {一次性写入提示的 K/V};
    \draw[<->, warn, line width=0.7pt] (2,0) -- (8,0);      \node[note, anchor=west] at (9,0) {每步读取全部权重与缓存};
  \end{scope}
\end{tikzpicture}
\caption{一次推理请求的两个阶段。预填充并行处理全部输入、写入 KV 缓存并给出首个 token；之后每个解码步只生成一个 token，却要读取全部权重与已有缓存}
\label{fig:prefill}
\end{figure}


两个阶段的特性不同，优化方向也不同：输入越长，预填充越耗时，TTFT 越大；输出越长，解码步数越多，总生成时间越长。

\subsection{KV 缓存：用存储换计算}
\label{subsec:kvcache}

进入解码阶段后，每个新 token 都要与之前所有 token 的 Key 和 Value 做注意力运算。如果每生成一个 token 都把全部历史 token 的 Key 和 Value 重新算一遍，总计算量将随序列长度呈平方增长。这种重复能否避免？

注意到历史 token 的 Key 和 Value 一旦算出就不会改变，因此可以把它们保存下来，每步只为新 token 计算一次并追加到缓存中。这就是 \term{KV 缓存}（\en{KV cache}）。需要强调的是，KV 缓存保存的是注意力计算所需的中间状态，而不是已经生成好的答案。

省下的计算要用显存来换。KV 缓存的大小可估算为
\begin{equation}
  M_{\mathrm{KV}} \approx 2 \times L \times H_{\mathrm{KV}} \times d_{\mathrm{head}} \times T \times B \times b_{\mathrm{KV}},
  \label{eq:kv}
\end{equation}
其中系数 2 对应 Key 和 Value，$L$ 为层数，$H_{\mathrm{KV}}$ 为 KV 头数，$d_{\mathrm{head}}$ 为每个头的维度，$T$ 为序列长度，$B$ 为同时驻留的请求数，$b_{\mathrm{KV}}$ 为每个元素的字节数。

\begin{example}[KV 缓存与并发上限]{ex:kv}
某 8B 模型有 32 层，采用分组查询注意力（GQA），KV 头数为 8，每头 128 维，以 BF16 存储（2 字节）。单个请求、序列长度 4096 时，
\[
  M_{\mathrm{KV}} = 2\times32\times8\times128\times4096\times1\times2 \approx 5.4\times10^{8}\ \text{字节} = \SI{0.5}{GiB}.
\]
在 \SI{80}{GB} 显卡上，BF16 权重约占 \SI{16}{GB}。即使不计其他开销，剩余显存也只能容纳约 $64/0.54\approx 119$ 个这样的请求。若序列长度增加到 32K，单个请求的 KV 缓存将达到 \SI{4}{GiB}，并发上限随之降至约 15 个。
\end{example}

可见，KV 缓存随序列长度和并发数线性增长，往往比权重本身更早耗尽显存（图\,\ref{fig:kv}），成为长上下文、高并发服务的主要约束。计算时要注意：采用 GQA 或 MQA（多查询注意力）的模型，$H_{\mathrm{KV}}$ 小于查询头数，不能误用查询头数代入。

\begin{figure}[htbp]
\centering
\begin{tikzpicture}[fig]
\begin{axis}[
  width=0.8\linewidth, height=5.6cm,
  xmin=0, xmax=32, ymin=0, ymax=100,
  xlabel={序列长度 T / K token}, ylabel={KV 缓存 / GiB},
  label style={font=\figfont\small}, tick label style={font=\figfont\small},
  xtick={0,4,8,16,24,32}, ytick={0,16,32,64,96},
  axis line style={figline}, tick style={figline},
  restrict y to domain=0:100, clip=true,
]
  \addplot[draw=rulegray, line width=1pt, domain=0:32, samples=2] {0.125*x};
  \addplot[draw=figline, line width=1pt, domain=0:32, samples=2] {2*x};
  \addplot[draw=accent, line width=1.2pt, domain=0:12.5, samples=2] {8*x};
  \addplot[draw=warn, dashed, line width=0.8pt, domain=0:32, samples=2] {64};
  \node[note, anchor=south east, text=warn] at (axis cs:31.5,64) {扣除权重后的可用显存 ≈ 64 GB};
  \node[note, anchor=east, text=accent] at (axis cs:8.2,74) {B = 64};
  \node[note, anchor=north west] at (axis cs:22,40) {B = 16};
  \node[note, anchor=south east] at (axis cs:31,5) {B = 1};
\end{axis}
\end{tikzpicture}
\caption{例\,\ref{ex:kv} 中 8B 模型的 KV 缓存随序列长度 $T$ 与并发请求数 $B$ 的变化（BF16）。曲线与虚线的交点给出了该配置下能够支撑的最大上下文长度}
\label{fig:kv}
\end{figure}


\subsection{推理优化技术}
\label{subsec:infer-opt}

围绕上述瓶颈，形成了一系列推理优化技术。理解它们的关键，是分清每一种技术\emph{作用于哪个环节、解决什么问题}。

\paragraph{FlashAttention：减少注意力计算中的显存读写} 标准注意力实现会把 $T\times T$ 的注意力分数矩阵完整写入显存再读回，序列越长，访存开销越大。FlashAttention\cite{dao2022flashattention}把计算分块，使中间结果尽量留在片上 SRAM 中，从而大幅减少对 HBM 的读写。它是\emph{精确}的注意力计算，结果与标准实现在数学上等价，并非近似算法；但不同的输入形状和硬件所获得的加速效果并不相同。

\paragraph{PagedAttention：像管理内存页一样管理 KV 缓存} 不同请求的输出长度事先未知，若按最大长度为每个请求预留连续显存，就会造成大量浪费和碎片。vLLM 提出的 PagedAttention\cite{kwon2023vllm}借鉴操作系统的分页虚拟内存：把 KV 缓存切成固定大小的块，按需分配，逻辑上连续、物理上可以分散。这显著提高了显存利用率，使同一张卡能容纳更多并发请求。

\paragraph{连续批处理：在每个解码步动态调度请求} 传统的静态批处理要等一批请求全部完成后，才能开始下一批，先结束的请求只能空等。\term{连续批处理}（\en{continuous batching}）\cite{yu2022orca}在每个解码步检查请求状态：已完成的请求立即移出，新到达的请求随即加入。这使 GPU 始终保持较大的有效批大小，从而提高总吞吐。

\paragraph{量化：用更少的位数表示数值} \term{量化}用低精度格式存储权重（有时也包括激活值和 KV 缓存）：FP16 每参数 2 字节，INT8 为 1 字节，4 位量化约为 0.5 字节。由于解码阶段访存受限，减少权重字节数通常有助于提高输出速度，同时降低显存占用。代价是需要额外存储缩放系数等元数据、进行反量化计算，并可能导致精度下降，因此应在目标任务的数据上验证效果，而不能只看速度\cite{dettmers2023qlora}。

表\,\ref{tab:infer-opt} 总结了这四种技术的作用位置。它们针对的瓶颈不同，可以组合使用，但不能相互替代。

\begin{table}[htbp]
\centering
\caption{常见推理优化技术的作用位置}
\label{tab:infer-opt}
\small
\begin{tabularx}{\linewidth}{@{}l l Y@{}}
\toprule
\heiti 技术 & \heiti 作用位置 & \heiti 主要解决的问题 \\
\midrule
FlashAttention & 注意力算子 & 注意力计算中对 HBM 的大量中间读写 \\
PagedAttention & KV 缓存的显存管理 & 预留连续显存造成的浪费和碎片，并发数受限 \\
连续批处理 & 请求调度 & 静态批次中先完成的请求空等，GPU 利用率低 \\
量化 & 数值表示 & 权重占用显存大，解码时访存量大 \\
\bottomrule
\end{tabularx}
\end{table}

\subsection{如何衡量推理性能}
\label{subsec:metrics}

性能指标必须对应具体的用户体验或成本目标。常用指标包括：
\begin{itemize}
  \item \emph{TTFT}：用户等待第一个字出现的时间，主要受排队和预填充影响；
  \item \emph{TPOT}（或输出速度）：连续生成时每个 token 的间隔，主要受解码阶段权重与 KV 缓存的读取影响；
  \item \emph{吞吐}：整个服务每秒处理的 token 数或完成的请求数，它与单个请求的速度是不同的指标；
  \item \emph{尾延迟}（如 P95、P99）：最慢的 5\% 或 1\% 请求的延迟。高并发时，部分请求可能远慢于平均水平，尾延迟往往比平均值更能反映真实体验。
\end{itemize}

这些指标之间存在权衡。增大批大小能让更多 token 共享权重读取，提高总吞吐；但也会占用更多 KV 缓存，并可能增加单个请求的排队时间和 TPOT。不存在能在所有负载下同时使 TTFT、吞吐和成本最优的调度策略。

\begin{pitfall}
报告“模型 A 比模型 B 快”时，必须同时说明硬件型号、精度、推理引擎及其版本、并发数，以及输入和输出长度的分布。缺少这些条件，性能数字既无法复现，也无法比较。
\end{pitfall}

\begin{keypoints}
\begin{itemize}[leftmargin=1.5em]
  \item 预填充计算受限，决定 TTFT；解码访存受限，决定输出速度。
  \item KV 缓存用显存换计算，随序列长度和并发数线性增长，可按式\,\eqref{eq:kv} 估算。
  \item FlashAttention、PagedAttention、连续批处理和量化分别作用于算子、显存管理、调度和数值表示。
\end{itemize}
\end{keypoints}
